% problem-000412 5 反钙钛矿离子导体
\section*{5. 反钙钛矿离子导体}
$
\mathbf {M} + \mathbf {M X A} + \mathbf {M Y} \rightarrow \mathbf {M} _ {3} \mathbf {X Y} + 1 / 2 \mathrm{A} _ {2} (\mathrm{g})
$

$\mathbf { A } _ { 2 }$ ⽆⾊⽆味。反应消耗 0.93 g M 可获得 0.50 L $\mathbf { A } _ { 2 }$ ⽓体 $( 2 5 ~ ^ { \circ } \mathrm { C } , 1 0 0 \mathrm { k P a } )$ 。（⽓体常量 $R = 8 . 3 1 4 \mathrm { { k P a } } \dot { . }$ L$\mathrm { m o l ^ { - 1 } K ^ { - 1 } } )$

\subsection*{5-1}
计算M的摩尔质量。

\subsection*{5-2}
A、M、X、Y各是什么？

\subsection*{5-3}
写出 $\mathbf { M } _ { 3 } \mathbf { X } \mathbf { Y }$ 发⽣⽔解的⽅程式。

\subsection*{5-4}
$\mathbf { M } _ { 3 } \mathbf { X } \mathbf { Y }$ 晶体属于⽴⽅晶系，若以X为正当晶胞的顶点，写出M和Y的坐标以及该晶体的最⼩重复单位。第6题（10分）某同学从书上得知，⼀定浓度的 $\mathrm { F e ^ { 2 + } }$ 、 $\mathrm { C r ( O H ) _ { 4 } } ^ { - } ,$ $\mathrm { N i ^ { 2 + } }$ 、 $\mathrm { M n O _ { 4 } }$ 2−和 $\mathrm { C u C l _ { 3 } }$ −的⽔溶液都呈绿$\mathtt { \ E q } _ { \mathtt { o } }$ 于是，请⽼师配制了这些离⼦的溶液。⽼师要求该同学⽤蒸馏⽔、稀硫酸以及试管、胶头滴管、⽩⾊点滴板等物品和尽可能少的步骤鉴别它们，从⽽了解这些离⼦溶液的颜⾊。请为该同学设计⼀个鉴别⽅案，⽤离⼦⽅程式表述反应并说明发⽣的现象（若A与B混合，必须写清是将A滴加到B中还是将B滴加到A中）。
